\scanpage{331}

is canonically endowed (not with a discrete \inserted{oriented} quasi-rigidification, which depends on the choice of an isomorphism $\overline{\mathbb Q}\simeq\overline{\mathbb Q}_0$ i.e. of an embedding $\overline{\mathbb Q}\subset\mathbb C$) but at least with a strict geometric Teichmüller group $\Sigma\subset\widehat{\widehat{\mathcal T}}(\pi)$, namely the image of $\pi_1(M_{g,\nu\overline{\mathbb Q}})$, \inserted{isomorphic to} This being so, We have a canonical homomorphism of group extensions

\diagram{(26)\quad
\begin{tikzcd}[ampersand replacement=\&, column sep=small]
1 \arrow[r] \&
\pi_1(M_{g,\nu\overline{\mathbb Q}}) \arrow[r] \arrow[d, "{\unsure{\wr|2}}"] \&
\pi_1(M_{g,\nu}) \arrow[r, no head] \arrow[d] \&
\operatorname{Gal}(\overline{\mathbb Q}/\mathbb Q) \arrow[r] \arrow[d] \&
1 \\
1 \arrow[r] \&
\Sigma \arrow[r] \&
\unsure{\mathcal N|N}_\Sigma \arrow[r] \arrow[d, equal] \&
\Gamma_\Sigma \arrow[r] \arrow[dr] \&
1 \\
\& \&
\operatorname{Norm}_{\widehat{\widehat{\mathcal T}}(\pi)}(\Sigma) \&
\&
\text{\margin{group of \deleted{\unsure{``arithmetic}}\inserted{outer} ``arithmetic'' automorphisms of $\pi$ (endowed with $\Sigma$)}}
\end{tikzcd}}

For the choice of an isomorphism $\overline{\mathbb Q}\simeq\overline{\mathbb Q}_0$ i.e. of an embedding $\overline{\mathbb Q}\subset\mathbb C$, the image of complex conjugation in $\operatorname{Gal}(\overline{\mathbb Q}_0/\mathbb Q)$ is none other than \unsure{the element} of order $2$ of $\Gamma_\Sigma$ associated with the geometric Teichmüller subgroup, associated with the canonical quasi-discretification (oriented, but no matter) on $\pi$ (defined by $\overline{\mathbb Q}\unsure{\simeq|\xrightarrow{\sim}}\overline{\mathbb Q}_0$).

\scanpage{332}
\mspage{168}

The natural conjecture here is that
\[
(27)\qquad \operatorname{Gal}(\overline{\mathbb Q}/\mathbb Q)\longrightarrow\Gamma_\Sigma
\]
from the Galois group to the arithmetic outer automorphisms of $(\pi,\Sigma)$ is an isomorphism. We can supplement it by the conjecture, still bolder, that moreover $\Sigma$ is invariant in $\widehat{\widehat{\mathcal T}}(\pi)$ ($\Longleftrightarrow\Sigma_{g,\nu}$ invariant in $\widehat{\widehat{\mathcal T}}_{g,\nu}$), which would thus mean that $\unsure{\mathcal N|N}_\Sigma=\widehat{\widehat{\mathcal T}}$, or again that \deleted{$\unsure{\pi_1(M_{g,\nu})}$}

\margin{Conjectural!!}
\[
(28)\qquad \pi_1(M_{g,\nu})\xrightarrow{\sim}\widehat{\widehat{\mathcal T}}(\pi)
=\operatorname{Aut\,ext}_{\mathrm{lac}}(\pi).
\]
and $\operatorname{Gal}(\overline{\mathbb Q}/\mathbb Q)$ would be identified with the ``arithmetic'' quotient of the group of outer automorphisms with loops of $\pi$.

Let now $U$ be an algebraic curve of type $g,\nu$ over $\overline{\mathbb Q}$. It is defined by a point of $M_{g,\nu\overline{\mathbb Q}}$ over $\operatorname{Spec}\overline{\mathbb Q}$, and \deleted{choosing a} as the fiber of $\mathcal U_{g,\nu}$ at this point. Let us choose a universal covering of $\unsure{U}$, whence a group $\pi(U)$, canonically isomorphic to the group $\pi$ a-

\scanpage{333}

bove, once an isomorphism has been chosen between $\widetilde{\mathcal U}_{g,\nu}$ and the universal covering of $\mathcal U_{g,\nu}$ deduced from $\widetilde U$, which gives an indeterminacy in $\pi_1(\mathcal U_{g,\nu\overline{\mathbb Q}})\simeq\unsure{\widehat{\mathcal S}^{+}}$.

Thus, once an isomorphism $\overline{\mathbb Q}\simeq\overline{\mathbb Q}_0$ has been chosen, one finds on $\pi_1(U)$ an oriented quasi-discretification (obvious moreover to define \unsure{directly}, by \unsure{transcendental} means) changed into its opposite by changing the isomorphism $\unsure{\overline{\mathbb Q}_2|\overline{\mathbb Q}}\simeq\overline{\mathbb Q}_0$ by complex conjugation. On the other hand, the subgroup $\widehat{\mathcal T}^{+}$ of $\widehat{\widehat{\mathcal T}}(\pi)$ corresponding to this quasi-discretification (oriented, little do we care) does not depend on the choice of the isomorphism in question $\overline{\mathbb Q}\simeq\overline{\mathbb Q}_0$, it is canonically associated with $\pi$ as the fundamental group of a $U$ over \deleted{$\unsure{\overline{\mathbb Q}}$} an algebraically closed field

[\uline{N.B.} one could take any algebraically closed field of characteristic $0$, no need for it to be \inserted{\unsure{a|over}} $\overline{\mathbb Q}$, cf. below].

\scanpage{334}
\mspage{169}

\deleted{Let us now consider an algebraic curve $U$ of type $(g,\nu)$ defined over a field $K$ of finite type --- thus defined by a morphism of $\mathbb Q$-multiplicities
\[
\operatorname{Spec}K\mathrel{\text{---}}M_{g,\nu}
\]
as the inverse image of $\mathcal U_{g,\nu}$. We thus canonically have a cartesian \unsure{homomorphism} fundamental groupoids
\diagram{\begin{tikzcd}[ampersand replacement=\&]
\Pi_1(U) \arrow[r] \arrow[d] \& \Pi_1(\mathcal U_{g,\nu}) \arrow[d] \\
\Pi_1(K) \arrow[r,dash] \& \Pi_1(M_{g,\nu})
\end{tikzcd}}
}

If the conjecture on $\Gamma$ holds, it would follow that the datum of a \deleted{\unsure{outer}} profinite group with loops \inserted{$\pi$} of type $(g,\nu)$, endowed with a strict geometric Teichmüller group $\Sigma\subset$ $\widehat{\widehat{\mathcal T}}(\pi)$ (which perhaps is not an additional structure at all --- $\Sigma$ would be uniquely determined by $\pi$ !) : would be equivalent to the datum $1^{\circ})$ of an \unsure{algebraically} closed extension $\overline{\mathbb Q}$ of $\mathbb Q$ (whose Galois group would be $N_{\widehat{\widehat{\mathcal T}}}(\Sigma)/\Sigma=\Gamma_\Sigma$, and the \inserted{\unsure{$\Gamma_0$-|$\Pi_0$-}}torsor $\operatorname{Isom}(\overline{\mathbb Q}_0,\overline{\mathbb Q}_1)$ would be identified

\scanpage{335}

with the set \inserted{($\subset\operatorname{Isom}(\widehat{\Pi}_{g,\nu},\pi)/\widehat{\mathcal T}_{g,\nu}^{\,0}$} of \deleted{\illegible{1}} \uline{oriented} quasi-rigidifications giving rise to \unsure{$\pi_{g,\nu}$|$\Pi_{g,\nu}$}) and of a universal covering of $\mathcal U_{g,\nu\,\overline{\mathbb Q}}$ --- or \unsure{simply}, of a universal covering $\widetilde{\mathcal U}_{g,\nu}$ of $\mathcal U_{g,\nu}=\mathcal U_{g,\nu\,\mathbb Q}$.

\margin{If one gives oneself \unsure{only} $\pi$ as an outer group, with $\Sigma$, this amounts to the datum of a universal covering of $M_{g,\nu}$.}

The datum of an isomorphism $\overline{\mathbb Q}\simeq\overline{\mathbb Q}_0$ thus amounts to that of an oriented quasi-rigidification on $\pi$ compatible \unsure{with} $\Sigma$ (condition perhaps automatically satisfied\ldots).
\margin{\uline{N.B.} If one gives up choosing $\widetilde U$, this \unsure{is equivalent|amounts} to working \unsure{with} \uline{outer} groups with loops.}
That said, the datum of an algebraic curve \inserted{of type $g,\nu$} over \inserted{\deleted{and of a universal covering}} \deleted{\unsure{(}\illegible{1} \unsure{would correspond}} an algebraically closed field $\overline{\mathbb Q}$, and of a universal covering \deleted{of it} of it would amount to the datum of a preceding datum (\unsure{namely} \inserted{\deleted{a \unsure{profinite outer group}} a universal covering of an $M_{g,\nu}$}), \uline{plus} a point of $M_{g,\nu,\overline{\mathbb Q}}$ over $\overline{\mathbb Q}$ \deleted{\unsure{this}} \inserted{\unsure{defining this universal covering} --- and,} this last datum would be identified with a \deleted{\unsure{point}} germ of splitting of the extension
\inserted{\[
1\longrightarrow\pi_1(M_{g,\nu\,\overline{\mathbb Q}})\longrightarrow\pi_1(M_{g,\nu})\longrightarrow\pi_1(\mathbb Q)\longrightarrow1
\]}
\unsure{The} \inserted{profinite} interpretation \inserted{(in terms of the \uline{outer} \unsure{groups} with loops $\pi$)} \inserted{\unsure{is then}} \deleted{\unsure{$\longrightarrow N_\Sigma\longrightarrow\Gamma_\Sigma\longrightarrow1$}} a $(\pi,\Sigma)$, and a germ of splitting of
\[
1\longrightarrow\Sigma\longrightarrow N_\Sigma\longrightarrow\Gamma_\Sigma\longrightarrow1.
\]
If one wants a curve over a finite subextension \unsure{$K'$} of $\overline{\mathbb Q}/\mathbb Q$, corresponding to a subgroup of finite index $\Gamma'\subset\Gamma_\Sigma$, it will be a matter of a partial \deleted{\illegible{1}} splitting $\Gamma'\longrightarrow N_\Sigma$.

\scanpage{336}
\mspage{170}

But surely, if this \deleted{\unsure{representation}} \inserted{description} of anabelian algebraic curves is fully faithful, it is not $2$-faithful, i.e. conditions on a splitting are needed, which will have to be examined later.

The description of an \deleted{\unsure{group}} algebraic curve of type $g,\nu$ over an arbitrary \inserted{fixed} field of finite type \inserted{\unsure{, characteristic\ $0$|over $\mathbb Q$}} $K$, endowed with an algebraically closed extension $\overline K$ (thus \unsure{$\Gamma=\operatorname{Gal}\overline K/K'$} $\longrightarrow\Gamma_0=\operatorname{Gal}(\overline{\mathbb Q}/\mathbb Q)$, where $\overline{\mathbb Q}$ is the algebraic closure of $\mathbb Q$ in $\overline K$), in terms of $\Gamma\longrightarrow\Gamma_0$, would be the following: datum of an outer group \unsure{with} loops $\pi$ of type $g,\nu$, of a Teichmüller \unsure{subgroup} $\Sigma$ in $\widehat{\widehat{\mathcal T}}(\pi)$, of an isomorphism $\overline{\mathbb Q}\simeq$ field defined by this situation (having as Galois group \deleted{by} $N_\Sigma/\Sigma$, so that we find an isomorphism $\Gamma_0\simeq N_\Sigma/\Sigma$) finally of a lifting of $\Gamma\longrightarrow\Gamma_0$ \deleted{\illegible{2}} \inserted{(which describes} an arithmetically outer action of $\Gamma$ on $(\pi,\Sigma)$) into an outer action $\Gamma\longrightarrow N_\Sigma$ of $\Gamma$ on $(\pi,\Sigma)$ ; with conditions which one will have to \unsure{try}

\scanpage{337}
to bring out (necessary and sufficient conjecturally) on this lifting.

\deleted{In order that}
Let us note that every embedding $\overline K\subset\mathbb C$
defines \inserted{\unsure{canonically|canonically}} on the outer group
$\pi=\pi_1(U_{\overline K})$
an oriented quasi-discretification, \unsure{replaced}
by the opposite when the embedding is replaced by its complex conjugate.
\deleted{Let us suppose} \deleted{\unsure{that|that}} I say that these
quasi-discretifications all define $\Sigma$ as the group of
\unsure{automorphisms}, and that the quasi-discretification
defined by $\overline K\hookrightarrow\mathbb C$ depends only on
its restriction to
$\overline{\mathbb Q}\subset\mathbb C$
--- in a more precise \unsure{way}, it is the one one defines
\unsure{directly} with the help of $\overline{\mathbb Q}\subset\mathbb C$ i.e. of
$\overline{\mathbb Q}\simeq\overline{\mathbb Q}_0$
--- But there is reason also to examine the way to obtain
\uline{discretifications}. Thus, if \inserted{\unsure{one has a|a}} $U$ is defined
over $\overline{\mathbb Q}_0$ (i.e. an outer \unsure{$\pi$|$\Pi$} of type
$(g,\nu)$, endowed with an oriented quasi-discretification

\scanpage{338}
\mspage{171}
\unsure{$\pi_0^{\natural+}$|$\pi_0^{\zeta+}$},
and with a germ of \deleted{\unsure{liftings}}\inserted{splitting}
of
\diagram{
\begin{tikzcd}[ampersand replacement=\&, column sep=large]
1 \arrow[r] \&
\Sigma_{\deleted{\unsure{\pi_0^{\natural+}}}}
\arrow[r] \arrow[d, equal] \&
N_{\Sigma} \arrow[r] \& \Gamma_{\Sigma} \arrow[r] \& 1 \\
\& \mathrm{Aut\,ext}_{\unsure{\mathrm{lac}}}
 (\pi,\unsure{\pi_0^{\natural+}|\pi_0^{\zeta+}}) \& \& \&
\end{tikzcd}
}
\unsure{)|,}
which makes the kernel \unsure{of|of the} group
defined by $\Gamma_{\Sigma}$ operate \inserted{outerly} on $\pi$, the rational points/$\overline{\mathbb Q}_0$
\unsure{corresponding|correspond} to the
\inserted{$\pi$-conjugacy classes of the} germs of liftings of
this operation\unsure{,} \unsure{outer|outer} into a genuine operation.
Let us therefore suppose such a point given, i.e.
\deleted{\unsure{$\Gamma_{\Sigma}$ operates}} \unsure{one has} an open subgroup
\unsure{$\Gamma^{\pi+}$|$\Gamma'$}$\subset\Gamma_{\Sigma}$
which does indeed operate on $\pi$,
\inserted{the operation defined} modulo inner automorphisms
(itself \unsure{unique}
since \unsure{$\pi^{\Gamma}$|$\overline\pi^{\Gamma}$}$=(1)$ !)
In this situation, it would be necessary to define in $\pi$
\inserted{$(\simeq\pi_1(U_{\unsure{P|p}},P))$}
a \uline{discretification}
\inserted{\uline{oriented} $\pi_0\subset\pi$,}
and not just an oriented quasi-discretification !
\deleted{But} (Of course, it must be in the class \unsure{which we} \inserted{have for ourselves} given in advance
of \deleted{quasi} oriented discretifications, \unsure{which we had}
precisely denoted \unsure{$\pi_0^{\natural+}$|$\pi_0^{\zeta+}$}.
We shall return to this --- and to the analogous situation over an arbitrary
base field $K$ of finite type\unsure{....|...}

\scanpage{339}

I \unsure{refer} to the beginning of the following \S\ (\S\,27) for the
change in terminology, ``quasi-discretification''
becoming a ``prediscretification''. But there is reason
also to introduce a finer notion, corresponding
--- \unsure{to the} \illegible{1} $\pi_1$ (profinite) of an algebraic \unsure{variety}
$X$ defined over an algebraically closed field $\overline K$, when one embeds
$\overline K$ into $\mathbb C$ : there is a discretification \unsure{modulo}
inner \unsure{automorphisms|automorphisms} --- we shall call it a
strict prediscretification. In the case of a
$\pi$ with loops, the homogeneous \deleted{space} \deleted{of} in $\widehat{\widehat{\mathcal T}}(\pi)$ of this is identified with
\[
\operatorname{Isom}_{\mathrm{lac}}(\widehat{\Pi}_{g,\nu},\pi)/\widehat{\Pi}_{g,\nu}\cdot
\overset{\deleted{\unsure{\widehat{\wedge}}}}{\mathcal T_{g,\nu}}
\inserted{\simeq\operatorname{Isom\,ext}_{\mathrm{lac}}(\widehat{\Pi}_{g,\nu},\pi)/\mathcal T_{g,\nu}}
\]
--- this notion \inserted{\unsure{also|also}} depends only on the outer \unsure{profinite} group with loops
defined by $\Pi_{g,\nu}$), and \deleted{\illegible{1}} \unsure{is equivalent}
to that of an ``outer basis'' in $\pi$
\inserted{\unsure{modulo the action of} $\mathcal T_{g,\nu}$,}
\inserted{i.e. of an isomorphism}
of $\pi$ with the ``standard'' $\widehat{\Pi}_{g,\nu}$ ;
\inserted{\unsure{modulo the action of} $\mathcal T_{g,\nu}$.}
The map
\[
\overset{\inserted{\text{Outer bases of }\pi}}{\deleted{\text{Strict prediscretifications }(\pi)}}
\longrightarrow \text{Prediscretifications of }(\pi)
\]
makes \unsure{therefore} of the left-hand side a relative torsor
\inserted{on the right,} over the one on the right, under the group $\mathcal T_{g,\nu}$.
\deleted{\illegible{1}}

A prediscretification of $(\pi,\Sigma)$ corresponds to a
choice of an isomorphism of the algebraic closure $\overline{\mathbb Q}_{\unsure{(\pi,\Sigma)}}$ \inserted{$=\overline{\mathbb Q}$} \unsure{associated}
with $(\pi,\Sigma)$, with $\overline{\mathbb Q}_0$. When one has given $(\pi,\Sigma)$ as
\unsure{corresponding} to an algebraic curve over $\overline{\mathbb Q}$, i.e. that one
\unsure{gives oneself|has given oneself} a \unsure{germ} of liftings of $\Gamma_\Sigma$ into $N_\Sigma\subset\widehat{\widehat{\mathcal T}}(\pi)$
i.e. a \unsure{germ|group} of outer action of $\Gamma_\Sigma$ on $\pi$, \unsure{every|whose}
prediscretification \inserted{must} give rise to a strict
discretification --- and \inserted{a} \deleted{\illegible{1}} discretification \unsure{is equivalent} \unsure{on lifting back}
a \unsure{germ} of action (``admissible'') of $\Gamma_\Sigma$ on $\pi$ \unsure{\ldots}

\scanpage{340}
\mspage{172}

27. \uline{Change of type} $(g,\nu)$ : a) \uline{plugging holes}.

I am going to change terminology, calling
prediscretification what I had called (a little
pejoratively !) \unsure{quasi}-discretification. Indeed, we
\unsure{foresee} that a prediscretification \inserted{--- compatible with $\Sigma$ \unsure{---)}} (which amounts \unsure{morally|morally}
to the embedding \inserted{(modulo complex conjugation)} into $\mathbb C$ of the algebraic closure $\overline{\mathbb Q}$ of $\mathbb Q$
canonically associated \unsure{with|$*$} $(\pi,\Sigma(\subset\widehat{\widehat{\mathcal T}}(\pi)))$, gives
rise as soon as the outer-arithmetic action of
$\Gamma_\Sigma=\unsure{N_\Sigma/\Sigma}$ on $\pi$ is lifted to a germ of action
on $\pi$, to a genuine discretification of $\pi$. \unsure{Ditto|The same} for
oriented prediscretifications. The group $N_\Sigma/\Sigma$ deserves
\margin{\unsure{morally|normally}, $\Sigma=\unsure{\widehat{\mathcal T}}$ (but there is no $\unsure{\mathcal T}$ !)}
a name as well --- I am going to call it the \uuline{arithmetic}
\uline{Teichmüller} \uline{group} associated with $(\pi,\Sigma)$, and
its elements will be called the
arithmetically outer automorphisms of $\pi$ --- \unsure{deduced} from an
outer automorphism (normalizing $\Sigma$) : by disregarding the
``geometric'' outer automorphisms (i.e. precisely those
$\in\Sigma$) --- \deleted{\unsure{of}} as the \inserted{ordinary} outer automorphisms \unsure{were}
\unsure{described} by disregarding the inner automorphisms of $\pi$.
We \deleted{\illegible{1}} shall note that the ``multiplier'' character
\[
\chi:\widehat{\widehat{\mathcal T}}(\pi)\unsure{\longrightarrow}\widehat{\mathbb Z}^{*}
\]
is trivial on $\Sigma$, by construction of $\Sigma$, and
passes to $\Gamma_\Sigma$
\[
\chi_\Sigma:\Gamma_\Sigma\longrightarrow\widehat{\mathbb Z}^{*}
\]